% Draft 2 -- heterogeneous experts under one calibrated gate.
% Target: TMLR (switch to the official tmlr.sty template before submission).
% Numbers: only \res... macros (generated by scripts/make_tables.py) or \pending{...}.
\documentclass[11pt]{article}
\usepackage[margin=1in]{geometry}
\usepackage{amsmath,amssymb,amsthm,booktabs,graphicx,xcolor}
\usepackage[numbers]{natbib}
\usepackage[colorlinks=true,citecolor=blue,linkcolor=red,urlcolor=blue]{hyperref}

\newcommand{\pending}[1]{{\color{red}\textbf{[#1]}}}
\newcommand{\todo}[1]{{\color{orange}\textit{[TODO: #1]}}}
\newtheorem{remark}{Remark}
% Generated by scripts/make_tables.py -- do not edit by hand.
\newcommand{\resTrackACorrectnessCases}{4,044,460}
\newcommand{\resTrackACorrectnessErrors}{0}
\newcommand{\resTrackAAddParams}{11,112}
\newcommand{\resTrackASubParams}{17,870}
\newcommand{\resTrackAMulParams}{310,251}
\newcommand{\resTrackADivParams}{99,607}
\newcommand{\resTrackAParserParams}{193,609}
\newcommand{\resTrackAAddNonzero}{482}
\newcommand{\resTrackASubNonzero}{983}
\newcommand{\resTrackAMulNonzero}{7,349}
\newcommand{\resTrackADivNonzero}{4,206}
\newcommand{\resTrackAParserNonzero}{1,500}
\newcommand{\resTrackALenCompiledAddThirtyTwo}{100}
\newcommand{\resTrackALenCompiledSubThirtyTwo}{100}
\newcommand{\resTrackALenCompiledMulThirtyTwo}{100}
\newcommand{\resTrackALenCompiledDivThirtyTwo}{100}
\newcommand{\resTrackALenPythonAddThirtyTwo}{100}
\newcommand{\resTrackALenPythonSubThirtyTwo}{100}
\newcommand{\resTrackALenPythonMulThirtyTwo}{100}
\newcommand{\resTrackALenPythonDivThirtyTwo}{100}
\newcommand{\resTrackALenQwenHalfBAddThirtyTwo}{0}
\newcommand{\resTrackALenQwenHalfBSubThirtyTwo}{0}
\newcommand{\resTrackALenQwenHalfBMulThirtyTwo}{0}
\newcommand{\resTrackALenQwenHalfBDivThirtyTwo}{0}
\newcommand{\resTrackALenQwenOneHalfBAddThirtyTwo}{0}
\newcommand{\resTrackALenQwenOneHalfBSubThirtyTwo}{0}
\newcommand{\resTrackALenQwenOneHalfBMulThirtyTwo}{0}
\newcommand{\resTrackALenQwenOneHalfBDivThirtyTwo}{0}
\newcommand{\resTrackANonzeroPctMin}{0.8}
\newcommand{\resTrackANonzeroPctMax}{5.5}
\newcommand{\resTrackAAddTwoFiftySixTokens}{776}
\newcommand{\resTrackAAddTwoFiftySixSeconds}{0.14}
\newcommand{\resTrackAMulThirtyTwoTokens}{5,284}
\newcommand{\resTrackAMulThirtyTwoSeconds}{4.15}
\newcommand{\resTrackADivThirtyTwoTokens}{11,635}
\newcommand{\resTrackADivThirtyTwoSeconds}{21.15}
\newcommand{\resTrackAParserDepthFiveTokens}{5,803}
\newcommand{\resTrackAParserDepthFiveSeconds}{12.90}


\title{Different Data, Different Experts: Calibrated Token-Level Routing over Trained, Compiled, and Tool Experts}
\author{Farshid Mossaiby\\Department of Civil Engineering, University of Isfahan\\\texttt{mossaiby@eng.ui.ac.ir}}
\date{Draft 2 -- \today}

\begin{document}
\maketitle

\begin{abstract}
\todo{Write last. Skeleton:}
Language models are trained as one dense network on a mixture of factual text, procedural
material such as arithmetic and code, and general prose, although these data types call for
different training: facts should be stored verbatim and be removable, procedures should
generalise beyond the lengths seen in training, and prose needs distribution matching. We
study the alternative of training one expert per data type, each under its own regime, and
combining them with a learned per-token gate that can also abstain. Experts can be trained
networks, \emph{compiled} transformers whose weights are written by hand, or external tools,
and all enter through the same calibrated read~\citep{mossaiby_roas}. \pending{Track C
headline: experts + gate vs.\ monolithic model at matched compute.} \pending{Track B headline:
pretrained experts sharing a tokenizer; swap test.} Compiled arithmetic experts and an
expression-parser expert that calls them are exact on \resTrackACorrectnessCases{} test cases
with \resTrackACorrectnessErrors{} errors, at operand lengths where prompted models of up to
1.5B parameters fail. \pending{Tools: abstention and revocation results.}
\end{abstract}

% =====================================================================================
\section{Introduction}
\label{sec:intro}

A dense language model is asked to do three different things with the same parameters:
remember facts, execute procedures, and produce fluent text. These place opposing demands on
training. Facts are best stored verbatim and must be updatable and removable, while gradient
updates to a dense network spread them across shared weights and make removal
approximate~\citep{meng2022rome}. Arithmetic and other procedures must generalise beyond the
lengths seen in training, a property that memorisation actively harms~\citep{mcleish2024abacus}.
Prose and code need the broad distribution matching that standard pretraining provides.

We ask whether these demands are better met by \emph{separate} experts, each trained under the
regime its data type calls for, and composed at decoding time by a learned per-token gate. The
gate also has an explicit abstention outcome, the null atom of~\citet{mossaiby_roas}, so that an
expert speaks only when it is relevant. Because composition happens in the shared output space,
an expert need not be a trained network at all: it can be a \emph{compiled} transformer whose
weights are written by hand to implement an algorithm exactly, or an external tool.

We test three hypotheses:
\begin{description}
\item[H1 (specialisation).] At matched total training compute, per-type experts combined by the
gate perform at least as well as one monolithic model, per data type and on mixed text.
\item[H2 (modularity).] Replacing or upgrading one expert changes performance only on its own
data type, and the gate adapts with zero or few training steps.
\item[H3 (frozen backbone).] When the backbone cannot be changed (for example, weights fixed in
hardware), external atoms provide insertion, update, removal, and tool revocation at a measured
cost.
\end{description}

\paragraph{Contributions.} \todo{finalise after results}
(i) A per-token gate over heterogeneous experts---trained, compiled, and tools---with calibrated
abstention, building on the read operator of~\citet{mossaiby_roas}.
(ii) Compiled experts for exact decimal $+,-,\times,\div$ of any length and an expression-parser
expert that calls them through the same latching mechanism; no weight depends on operand length
(Section~\ref{sec:compiled}).
(iii) A controlled from-scratch study of H1 (Section~\ref{sec:scratch}).
(iv) Experiments with pretrained experts that share a tokenizer, including expert swaps
(Section~\ref{sec:pretrained}).
(v) Frozen-backbone fact removal and tool revocation (Section~\ref{sec:frozen}).

% =====================================================================================
\section{Related work}
\label{sec:related}

\paragraph{Separately trained experts.} Branch-Train-Merge~\citep{li2022btm} and its clustered
variant~\citep{gururangan2023cbtm} train domain experts independently and ensemble them;
DEMix~\citep{gururangan2022demix} uses domain-specific feed-forward layers; Branch-Train-MiX
(BTX)~\citep{sukhbaatar2024btx} trains math, code, and knowledge experts from a common seed and
merges them into a mixture-of-experts model. These works split data by \emph{domain} under one
training regime. We split by \emph{data type} with a different regime per type, admit experts
that are not trained at all, and add an abstention outcome. BTX is the closest baseline for H1.

\paragraph{Token-level collaboration between models.} Co-LLM~\citep{shen2024colm} learns, per
token, whether a base model defers to an expert; DExperts~\citep{liu2021dexperts} and
proxy-tuning~\citep{liu2024proxy} combine models in logit space; CALM~\citep{bansal2024calm}
composes models through cross-attention. Our gate is closest to Co-LLM; it differs in routing
over more than two sources of different kinds, in its calibrated null atom, and in latching
onto an expert for the duration of its answer.

\paragraph{Tools as output elements.} Toolformer~\citep{schick2023toolformer} and
ReAct~\citep{yao2023react} emit tool calls as plain text, while ToolkenGPT~\citep{hao2023toolkengpt} and
ToolGen~\citep{wang2024toolgen} represent each tool as a vocabulary token. Representing a tool as
an element of the output space is therefore not new. What we add is that tools sit in the same
calibrated read as memory and parameter experts: one gate across spaces, abstention when no tool
fits, revocation by anti-atoms, and fallback to stored records.

\paragraph{Compiled transformers.} RASP~\citep{weiss2021rasp} and Tracr~\citep{lindner2023tracr}
compile programs into transformer weights; looped transformers can emulate a programmable
computer~\citep{giannou2023looped}; transformers with chain of thought can simulate
polynomial-time computation~\citep{merrill2024cot}. Our compiled experts use a scratchpad with
fixed row layouts so that weights do not depend on input length, and they run \emph{inside} a
language model's decoding loop as experts.

\paragraph{Arithmetic in language models.} Position encodings designed for digits
(Abacus,~\citealp{mcleish2024abacus}), scratchpads, and program-aided
reasoning~\citep{gao2023pal} improve learned arithmetic. We compare a compiled expert against
prompted models, a trained arithmetic model, and a Python tool (Section~\ref{sec:compiled}).

\paragraph{Editing and unlearning.} \todo{ROME/MEMIT/GRACE/SERAC/WISE; TOFU~\citep{maini2024tofu};
the frozen-backbone setting; cite RoAS for anti-atom vs deletion.}

% =====================================================================================
\section{Method}
\label{sec:method}
\todo{Trim to essentials; cite RoAS for the read operator, null atoms, anti-atoms.}

\paragraph{Experts and the read.} \todo{Each expert $e$ produces next-token logits $\ell_e(x)$
over a shared vocabulary; the gate $\lambda(x)$ is a distribution over experts and a null atom.
Three combination rules: probability mixture, logit mixture (product of experts), hard routing.
State when each is used.}

\begin{remark}[Veto]
If an expert assigns $+C$ to one token and $-C$ to all others, and every other source has logits
bounded by $B$, the logit mixture selects that token whenever the expert's gate weight exceeds
$B/(B+C)$.
\end{remark}

\paragraph{Latching and call blocks.} \todo{Session latching, consumption guard; an expert may emit
a call block that the latch routes to another expert (Section~\ref{sec:parser}).}

\paragraph{Gate.} \todo{Inputs (last state, windowed mean and max), training without routing labels
by mixture likelihood, calibration of the null atom on held-out data as in RoAS.}

% =====================================================================================
\section{Compiled experts}
\label{sec:compiled}

A compiled expert is an ordinary decoder-only transformer---token embedding, a fixed
positional signal, residual blocks of causal softmax attention and ReLU MLPs, and an
unembedding---whose weights are written by hand. Nothing is trained. We build experts for
decimal addition, subtraction, multiplication, and division, and an expression-parser expert
that calls them.

\subsection{Three primitives}
\label{sec:primitives}
Positions are embedded as the pair $(p, p^2)$. This makes three operations exact:
\begin{itemize}
\item \textbf{Lookup.} A head whose score is $-\Lambda (p_k - t)^2$, written as a bilinear form in
$(p_k^2, p_k)$ and $(1, t)$, reads the token at exactly position $t$, provided $t$ is a linear
function of the residual stream. The correct key beats any other by a factor $e^{-\Lambda}$.
\item \textbf{Select.} A head whose key score is a linear function of key features attends to its
maximiser, for example the most recent row separator.
\item \textbf{Integer ReLU logic.} On integers, $[x\ge c] = \mathrm{ReLU}(x-c+1)-\mathrm{ReLU}(x-c)$
and $[x = c]$ is a three-ReLU triangle. Gates are integer expressions equal to $1$ when active and
$\le 0$ otherwise.
\end{itemize}
The design rule for every expert is: \emph{every position a step must read is a linear function
of the current position and of positions of marker tokens}. Under that rule no weight depends on
operand length.

\subsection{Addition and subtraction without a scratchpad}
Addition emits the most significant digit first, in one forward pass per digit. The carry into
digit $i$ is the carry flag of the first later digit pair whose sum is not $9$. It is found with
a linked list: each pair records whether its sum differs from $9$, each position records the
nearest such pair at or before it, and the first such pair after $i$ is the one whose
predecessor equals $i$'s. Subtraction is ten's complement through the same construction, with
an initial carry of $1$ read from the start token and a second chain for $b-a$. Both chains share
the linked list because $a_i + 9 - b_i \neq 9 \iff a_i \neq b_i$; the carry out of the first chain
is the sign.

\subsection{Multiplication and division with a scratchpad}
Multiplication writes one row per digit $b_r$ of the multiplier. Each row holds a copy of the
multiplier shifted by $r$ (so its last symbol is $b_r$), the multiplicand shifted by $r$, and the
running sum $S_r = S_{r-1} + A\,b_r\,10^r$, least significant first. A sum cell is a single token
whose embedding exposes both its digit and its carry, so the carry chain needs no extra tokens.
All rows have the same length $m + 2(n+m) + 1$, so every value a cell needs lies at a fixed
offset.

Division is restoring long division. Each row is a bring-down row or the $c$-th subtraction of
the current step, and carries the unconsumed dividend digits, the quotient so far, and the
working number with per-digit borrows. The next row's type depends only on the previous row:
after a bring-down comes the first subtraction; after a non-negative subtraction comes the next;
after a negative one comes a bring-down that appends the quotient digit $c-1$ and reuses the row
\emph{two} back, which is always the last non-negative value. Small state lives in token
identities and unbounded state is copied row to row at linear offsets.

\subsection{An expression parser that calls other experts}
\label{sec:parser}
The parser expert evaluates expressions with brackets, precedence, left associativity, and
unary minus. It rewrites the expression one step per row. For each binary operation it emits a
call block \texttt{[a op b?}; the latch routes \texttt{a op b} to the matching arithmetic expert
and appends its answer, which the parser copies into the next row. A group holding one number
loses its brackets. Finding the next operation needs no bracket counting: the first closing
bracket of a row closes an innermost group, its partner is the nearest opener before it, and
inside that group exactly one multiplicative (else additive) operator is marked as the first
since its opener. These facts are summarised at each row terminator, from which the writer of the
next row reads them.

\subsection{Results}
\label{sec:compiled_results}

\paragraph{Correctness.} Table~\ref{tab:correctness} lists every test suite: all operand pairs
below $10^2$ and below $10^3$ for each circuit (about four million cases), random operands with many 9s up to 200
digits (addition, subtraction) and 32 digits (multiplication, division), signed decimals through
the formatting layer, and random expressions and balanced trees for the parser, checked against
a reference evaluator with identical truncating-division semantics.
\begin{table}[t]
\centering\small
\caption{Correctness of the compiled experts. Every case is compared exactly with Python
integers, \texttt{Decimal}, or the reference expression evaluator.}
\label{tab:correctness}
\IfFileExists{tables/track_a_correctness.tex}{% Generated by scripts/make_tables.py -- do not edit by hand.
\begin{tabular}{@{}llrr@{}}
\toprule
Expert & Test suite & Cases & Errors \\
\midrule
add & exhaustive, operands $<$ $10^{2}$ & 10,000 & 0 \\
sub & exhaustive, operands $<$ $10^{2}$ & 10,000 & 0 \\
mul & exhaustive, operands $<$ $10^{2}$ & 10,000 & 0 \\
div & exhaustive, operands $<$ $10^{2}$ & 9,900 & 0 \\
add & exhaustive, operands $<$ $10^{3}$ & 1,000,000 & 0 \\
sub & exhaustive, operands $<$ $10^{3}$ & 1,000,000 & 0 \\
mul & exhaustive, operands $<$ $10^{3}$ & 1,000,000 & 0 \\
div & exhaustive, operands $<$ $10^{3}$ & 999,000 & 0 \\
add & random, up to 200 digits & 200 & 0 \\
sub & random, up to 200 digits & 200 & 0 \\
mul & random, up to 32 digits & 60 & 0 \\
div & random, up to 32 digits & 60 & 0 \\
add & signed decimals (arith) & 1,000 & 0 \\
sub & signed decimals (arith) & 1,000 & 0 \\
mul & signed decimals (arith) & 1,000 & 0 \\
div & signed decimals (arith) & 1,000 & 0 \\
parser & balanced trees, depth 1--4 & 40 & 0 \\
parser & random expressions & 1,000 & 0 \\
\bottomrule
\end{tabular}
}{\pending{table}}
\end{table}

\paragraph{Length generalisation.} Table~\ref{tab:length} compares exact-match accuracy against
operand length for the compiled experts, prompted Qwen2.5-Instruct models with 0.5B and 1.5B
parameters, and a Python tool. \pending{Add the trained arithmetic expert from
Section~\ref{sec:scratch}.}
\begin{table}[t]
\centering\small
\caption{Exact-match accuracy (\%) against the number of digits per operand (for division the
divisor has half as many). Same items for every system.}
\label{tab:length}
\IfFileExists{tables/track_a_length_generalisation.tex}{% Generated by scripts/make_tables.py -- do not edit by hand.
\begin{tabular}{@{}llrrrrrr@{}}
\toprule
Op & System & 1 & 2 & 4 & 8 & 16 & 32 \\
\midrule
add & compiled & 100 & 100 & 100 & 100 & 100 & 100 \\
 & Python & 100 & 100 & 100 & 100 & 100 & 100 \\
 & Qwen2.5-0.5B-Instruct & 97 & 80 & 50 & 3 & 0 & 0 \\
 & Qwen2.5-1.5B-Instruct & 100 & 97 & 87 & 50 & 3 & 0 \\
\midrule
sub & compiled & 100 & 100 & 100 & 100 & 100 & 100 \\
 & Python & 100 & 100 & 100 & 100 & 100 & 100 \\
 & Qwen2.5-0.5B-Instruct & 100 & 87 & 70 & 13 & 0 & 0 \\
 & Qwen2.5-1.5B-Instruct & 100 & 93 & 80 & 20 & 0 & 0 \\
\midrule
mul & compiled & 100 & 100 & 100 & 100 & 100 & 100 \\
 & Python & 100 & 100 & 100 & 100 & 100 & 100 \\
 & Qwen2.5-0.5B-Instruct & 97 & 27 & 0 & 0 & 0 & 0 \\
 & Qwen2.5-1.5B-Instruct & 97 & 53 & 0 & 0 & 0 & 0 \\
\midrule
div & compiled & 100 & 100 & 100 & 100 & 100 & 100 \\
 & Python & 100 & 100 & 100 & 100 & 100 & 100 \\
 & Qwen2.5-0.5B-Instruct & 80 & 17 & 0 & 0 & 0 & 0 \\
 & Qwen2.5-1.5B-Instruct & 73 & 97 & 30 & 0 & 0 & 0 \\
\bottomrule
\end{tabular}
}{\pending{table}}
\end{table}

\paragraph{Cost.} Table~\ref{tab:cost} gives the size of each expert. All are tiny next to a
language model, and sparse: only \resTrackANonzeroPctMin--\resTrackANonzeroPctMax\% of the
weights are nonzero, because every head reads and writes a few named features and every MLP unit
tests one integer condition. Size does not depend on operand length; length is paid in generated
tokens (Table~\ref{tab:scaling}). Addition and subtraction emit one token per digit
(\resTrackAAddTwoFiftySixTokens{} tokens in total for two 256-digit operands). The scratchpad for
multiplication grows as $m(2n+3m+1)$ (\resTrackAMulThirtyTwoTokens{} tokens for $32\times 32$
digits), and division grows as at most about $11n(2n+m+4)$ (\resTrackADivThirtyTwoTokens{} tokens
for 32 by 16 digits). The parser adds about one row per operation, and each row grows with the
expression (\resTrackAParserDepthFiveTokens{} tokens for a balanced tree with 31 operators).
Wall-clock times are for a single-threaded NumPy reference with a key/value cache. They are
\resTrackAAddTwoFiftySixSeconds\,s, \resTrackAMulThirtyTwoSeconds\,s and
\resTrackADivThirtyTwoSeconds\,s for these cases. The implementation is not optimised; the point
is exactness at every length, not speed. On an accelerator, the cost per token of an expert with
$d_{\mathrm{model}}\le 243$ is negligible next to that of the base model it is blended with, so
the real cost is the extra tokens, which Section~\ref{sec:ablations} measures.
\begin{table}[t]
\centering\small
\caption{Size of the compiled experts. ``Blocks'' counts attention and MLP sublayers. The last
column is the percentage of nonzero weights.}
\label{tab:cost}
\IfFileExists{tables/track_a_cost.tex}{% Generated by scripts/make_tables.py -- do not edit by hand.
\begin{tabular}{@{}lrrrrrrrr@{}}
\toprule
Expert & Blocks & Heads & MLP units & $d_{\mathrm{model}}$ & Vocab & Weights & Nonzero & \% \\
\midrule
add & 9 & 7 & 72 & 48 & 16 & 11,112 & 482 & 4.3 \\
sub & 9 & 8 & 132 & 49 & 17 & 17,870 & 983 & 5.5 \\
mul & 6 & 12 & 809 & 161 & 117 & 310,251 & 7,349 & 2.4 \\
div & 5 & 15 & 427 & 95 & 48 & 99,607 & 4,206 & 4.2 \\
parser & 17 & 28 & 181 & 243 & 26 & 193,609 & 1,500 & 0.8 \\
\bottomrule
\end{tabular}
}{\pending{table}}
\end{table}
\begin{table}[t]
\centering\small
\caption{Tokens per call (prompt and scratchpad) and wall-clock time of the NumPy reference
against operand size. The averages are over 5 random cases, or 2 cases above 16 digits.}
\label{tab:scaling}
\IfFileExists{tables/track_a_scaling.tex}{% Generated by scripts/make_tables.py -- do not edit by hand.
\begin{tabular}{@{}llrr@{}}
\toprule
Expert & Size & Tokens & Seconds \\
\midrule
add & 4 digits & 20 & 0.003 \\
add & 16 digits & 56 & 0.007 \\
add & 64 digits & 200 & 0.029 \\
add & 256 digits & 776 & 0.137 \\
sub & 4 digits & 20 & 0.003 \\
sub & 16 digits & 56 & 0.008 \\
sub & 64 digits & 200 & 0.031 \\
sub & 256 digits & 776 & 0.155 \\
mul & 2 digits & 34 & 0.008 \\
mul & 4 digits & 104 & 0.026 \\
mul & 8 digits & 364 & 0.100 \\
mul & 16 digits & 1,364 & 0.522 \\
mul & 32 digits & 5,284 & 4.151 \\
div & 4 digits & 273 & 0.079 \\
div & 8 digits & 831 & 0.358 \\
div & 16 digits & 2,967 & 1.947 \\
div & 32 digits & 11,635 & 21.149 \\
parser & depth 1 (1 calls) & 38 & 0.027 \\
parser & depth 2 (3 calls) & 117 & 0.069 \\
parser & depth 3 (7 calls) & 485 & 0.370 \\
parser & depth 4 (15 calls) & 1,578 & 1.598 \\
parser & depth 5 (31 calls) & 5,803 & 12.900 \\
\bottomrule
\end{tabular}
}{\pending{table}}
\end{table}

% =====================================================================================
\section{Training experts from scratch}
\label{sec:scratch}
\todo{Track C: data, regimes, monolithic baselines at matched compute, gate, results; go/no-go
outcome decides the framing of the paper.}

\section{Pretrained experts that share a tokenizer}
\label{sec:pretrained}
\todo{Track B: Qwen2.5 base/Coder/Math, cached outputs, staged build-up, swap tests.}

\section{Frozen backbone and tools}
\label{sec:frozen}
\todo{Track D (TOFU, MQuAKE) and D2 (API-Bank, BFCL irrelevance, revocation, fallback).}

\section{Ablations and cost}
\label{sec:ablations}
\todo{Gate inputs, window, latching, consumption guard, C and temperature, null atom, mixing
rule; FLOPs and latency.}

\section{Discussion and limitations}
\label{sec:discussion}
\todo{Write after results. Known limitations so far: compiled experts are slow in the reference
implementation and their scratchpads grow quadratically; float64 positions bound sequence length
to roughly $10^5$--$10^6$ tokens; repeated unary minus is not parsed; division truncates.}

\bibliographystyle{plainnat}
\bibliography{references_v2}
\end{document}
